% figure (2026-11-17): centre-of-mass kinematics of e- mu- -> e- mu- with m_e = 0, after PHY 513 Lecture 21 slide 11 and PS p. 154; generated by qft-C/tools/l21/figs/make_figs.sh
\tikzset{
  fdl/.style={line width=0.8pt, black},
  fdx/.style={circle, fill=black, inner sep=0pt, minimum size=3.6pt},
  fdv/.style={circle, fill=figB, draw=figB, inner sep=0pt, minimum size=5pt},
  fdlab/.style={font=\scriptsize, inner sep=2.6pt},
  fdmom/.style={->, >=Stealth, figO, line width=0.6pt},
}
\tikzset{
  fda/.style={fdl, postaction={decorate}, decoration={markings, mark=at position 0.55 with {\arrow{Stealth[length=5pt]}}}},
  fdS/.style={font=\small},
}
\tikzset{
  fline/.style={line width=0.8pt, black, postaction={decorate}, decoration={markings, mark=at position 0.56 with {\arrow{Stealth[length=5.5pt, width=4.5pt]}}}},
  sline/.style={line width=0.8pt, black, dash pattern=on 3.5pt off 2.2pt},
  blob/.style={draw, line width=0.8pt, fill=figB!12},
}
\tikzset{
  pline/.style={line width=0.8pt, black, decorate, decoration={snake, amplitude=1.6pt, segment length=5.2pt}},
}
\begin{tikzpicture}[font=\small, >=Stealth]
  \draw[->, line width=0.9pt] (-3.2,0)--(-0.2,0); \draw[->, line width=0.9pt] (3.2,0)--(0.2,0);
  \node[above] at (-2.0,0.05) {$e^-$: $p_1 = (k, k\hat{\mathbf z})$};
  \node[below] at (2.1,-0.05) {$\mu^-$: $p_2 = (E, -k\hat{\mathbf z})$};
  \draw[->, line width=0.9pt, figB] (0,0)--(1.9,1.55); \draw[->, line width=0.9pt, figB] (0,0)--(-1.9,-1.55);
  \node[figB, right] at (1.9,1.62) {$e^-$: $p_1' = (k, \mathbf k)$};
  \node[figB, left] at (-1.9,-1.62) {$\mu^-$: $p_2' = (E, -\mathbf k)$};
  \draw[gray] (0.95,0) arc (0:39.2:0.95); \node at (1.22,0.36) {$\theta$};
  \draw[->, gray] (3.6,0)--(4.3,0); \node[font=\scriptsize, gray, right] at (4.3,0) {$\hat{\mathbf z}$};
  \node[align=left, anchor=west, font=\footnotesize] at (5.4,0.2) {$|\mathbf k| = k$ \ (elastic; $m_e = 0$)\\[2pt] $E = \sqrt{k^2 + m_\mu^2}$\\[2pt] $\mathbf k\cdot\hat{\mathbf z} = k\cos\theta$\\[2pt] $E_{\rm cm} = E + k$};
\end{tikzpicture}
